\documentclass[11pt]{article}

\usepackage{a4wide}
\usepackage{latexsym}
\pagestyle{empty}

\def\CC{\bf C}

\title{About an algorithm of T. Oaku}

\author{Mar\'{\i}a Isabel Hartillo Hermoso}

\date{Universidad de C\'adiz}

\begin{document}

\maketitle

Let be $f\in \CC$$ [x]$ with $x=(x_1,\ldots,x_n)$, and ${\cal A}_n$ the Weyl
algebra with $n$ variables.
The polynomials
$b(s)\in\CC$$ [s]$ such that there exists $Q\in{\cal A}_{n}[s]$ with:
$$
Q(f^{s+1})=b(s)f^s,
$$
form an ideal of $\CC $$[s]$, which is a principal ideal domain.
The Bernstein polynomial
asociated to $f$ is this monic polynomial $b_f(s)$ of that ideal.

The proof given by Bernstein of the existence of $b_f(s)$ is not constructive,
and that is why there was not effective algorithm for the calculation of the
Bernstein polynomial.

In 1997, T. Oaku (``An algorithm of computing $b$-functions'', Duke
Mathematical Journal, {\bf 87}, vol. 1, 115--132) gave a method which calculate
the Bernstein polynomial. This algorithm used the computation of a Gr\"obner
basis of an ideal for an order $\prec_F$ that was not a well ordering. In
order to solve that problem, 
the author introduced the $F$-homogenization of the
operators, in the ring ${\cal A}_{n+1}[x_0]$. Then he defined a new order in
that ring, the homogenization order, that was a well ordering and such that 
there always exists
a $H$-Gr\"obner basis. Dehomogenizing that Gr\"obner basis he
returned to ${\cal A}_{n+1}$ and obtained a basis for the initial order
$\prec_F$. Using elimination tecniques he finally gave the 
Bernstein polynomial.

From this work, we deal with the problem as  a Gr\"obner basis computation
for a graded ideal in the graded algebra given by
Kashiwara--Malgrange filtration. Using a new method, we define now the
homogenization in the Rees algebra. This comes from ``Homogenising
differential operators'' by F. J.
Castro--Jim\'enez and L. Narv\'aez--Macarro, (Prepublicaciones de la Facultad
de Matem\'aticas, Universidad de Sevilla, {\bf 36}, 1998), permiting an ample
degree of freedom in the choice of the order. Let us introduce a new order
$\prec_V$, that order monomials of bigger $V$-order first and besides, 
is a elimination order in $x$ and $\partial$. This new method allows to
calculate just a Gr\"obner basis in the Rees algebra, (for a homogenization
order), and that basis chives directly to the Bernstein polynomial. 
With this new
algorithm we have a range of options for the elimination order. We finish
with the calculation of non trivial examples using the new method.

\vspace{.5cm}

Address of author:

Mar\'{\i}a Isabel Hartillo Hermoso

E. U. Empresariales
c/ Por--Vera, 54 

11403 
Jerez de la Frontera
(C\'adiz)

Phone: 956335762

Fax:956345104

E-mail: hartillo@deimos.us.es


\end{document}

